\section{Survey Methodology}
\label{sec:methodology}
\revise{This retrospective literature survey examines LLM-assisted security analysis, testing, verification, and repair across OS components, emphasizing approaches that ground their outputs in verifiable artifacts. Studies that protect LLM-based systems are discussed as adjacent work where relevant, rather than as direct evidence of improved OS security. The collected literature is not claimed to result from a complete, reproducible systematic search.}

\revise{For transparency, we use the search-reporting items of PRISMA-S}~\revcite{rethlefsen2021prismas} \revise{and the selection-reporting guidance of PRISMA 2020}~\revcite{page2021prisma} \revise{to document the available evidence and remaining gaps; referring to these guidelines does not imply that all their items have been satisfied.}

\paragraph{\textbf{Search strategy and venues}}
\revise{The original collection combined targeted searches and foundational references. The 18 surviving historical query formulations are in \url{reproduct/search/}; their database-specific forms, execution dates and results are not established. One formulation combines LLM, large language model or Transformer with operating system security or OS security. TEE and hardware-isolation terms occur, but coverage of neighboring VMM, hypercall, enclave-interface and secure-monitor terms is not established. Representative publication outlets include ICSE, USENIX Security, IEEE S\&P, ASPLOS, OSDI, NeurIPS and ICLR; these are outlets, not verified search databases. Preprints were also considered.}

\revise{On November 5, 2026, an exploratory update used Crossref after automated Google Scholar requests were blocked. It submitted 100 queries with a 2022-01-01 lower publication-date limit and retained up to 20 ranked results per query. Crossref received transformed free-text expressions, not the original Boolean queries. The queries, times, returned counts and status are recorded in \url{reproduct/search/search_metadata.csv}. This later update does not reconstruct the historical search.}

\revise{A separate audit on the same date saved up to 100 results for each of the 100 Crossref queries and for 12 arXiv queries. Crossref and arXiv used different query syntax and date limits; all Crossref queries and two arXiv queries were truncated. Exact DOI and arXiv-identifier matching grouped the 10,467 retrieved rows into 6,408 provisional records. The two truncated arXiv searches were subsequently paginated; their additional records and the time and ranking limitations are documented separately. These bounded searches are exploratory, and their results are not a verified replacement for the historical synthesis. Query logs, linkage rules and scripts are documented under \url{reproduct/search_protocol/} and \url{reproduct/arxiv_pagination/}; the research archive also retains provider responses.}

\revise{The audit records received preliminary triage of titles and available abstracts using \texttt{deepseek/deepseek-flash} via OpenCode. A separate call using the same model assessed the 2,311 records with available abstracts in the 6,408-record snapshot. These model outputs are not independent human screening or final eligibility decisions. Missing abstracts, mismatched evidence excerpts, disagreements and flagged OS-related records remain for review; a check of cited work found a potential model false exclusion. The screening protocol, decision fields, conservative rules and audit are documented under \url{reproduct/screening/} and \url{reproduct/audit/}; individual model inputs and responses are retained in the research archive.}

\revise{The associated search records, screening materials, inventories and supporting scripts are organized in the accompanying \texttt{reproduct/} directory. The exact contents of any public release depend on the materials made available with that release.}

\revise{The stated publication window for candidate research is 2020--2026; older foundational references are used for background rather than as evidence of recent LLM-assisted OS-security methods. Coverage in 2026 is limited to the collected records, not the entire calendar year. For the historical collection, the available materials do not establish exact search dates, a database-specific query log, or a consistent rule for choosing between first-publication and final-publication dates. The venues listed above identify publication outlets rather than a verified list of searched databases. These documentation gaps limit reproducibility and the completeness of the review.}


\paragraph{\textbf{Deduplication and publication versions}}
\revise{The following version-handling principles guide interpretation of the literature; the available spreadsheet does not document a completed record-to-study reconciliation.}
\revise{Version reconciliation requires matching DOI or arXiv identifiers and comparing normalized titles and author metadata; ambiguous or renamed records require comparison of the underlying studies. A preprint and its published counterpart should be treated as one study when they report the same contribution, with the published version preferred when available within the search period. Preprint-only work can be considered under the same substantive criteria, with its publication status identified. Distinct contributions in extended versions require explicit justification before separate counting. These principles do not imply that the 220 scored records already represent 220 unique studies.}

\paragraph{\textbf{Eligibility and screening}}
\revise{The criteria below define the intended scope of the core synthesis, rather than certifying the eligibility of every previously scored entry.}
For the core synthesis, a study must (i) fall within the stated date and document-type limits; (ii) use an LLM as a substantive part of a security analysis, testing, repair, or verification method; (iii) explicitly target an OS component or privileged interface; and (iv) provide sufficient methodological detail and evaluation evidence to assess that OS-security application. Evidence may include experiments, case studies, or formal results; public artifact availability is recorded rather than assumed to be mandatory.
Excluded from the core synthesis are duplicates, out-of-window studies, non-LLM methods, generic code-generation or application-security studies without an explicit OS target, and opinion or promotional material without assessable research evidence. Work that uses OS isolation or TEEs solely to protect LLM execution is treated as adjacent research rather than direct evidence that LLMs improve OS security. Surveys, foundational papers, standards, and vendor documentation may support background discussion but are not counted as independent core studies.
\revise{In particular, ERP, OA, and mail studies confined to application business logic or message-content analysis are excluded from the core synthesis.}
\revise{The surviving materials do not provide a complete record-level log of title/abstract decisions, full-text exclusions, or retrieval failures, nor do they establish whether language or full-text availability restrictions were applied consistently. Eligibility criteria describe the intended scope and do not certify a final eligible-study set.}

\paragraph{\textbf{\revise{Relevance scoring and evidence}}}
\revise{The exploratory update returned 2,000 Crossref result rows across the 100 queries. DOI-based merging yielded 1,454 candidate records; records without a DOI were kept separately. A ranked subset of 200 was inspected rapidly using titles and available abstracts, and 101 entries were selected as direct or adjacent candidates. Six matched previously retained records and 95 were added as new rows. The resulting inventory contains 130 retained legacy rows and 95 new rows, totaling 225. Earlier low-relevance records were not systematically carried forward. These are candidate records, not a census of unique eligible studies. Matching used recorded identifiers or normalized titles where available; unresolved duplicates and preprint/publication versions can remain. The candidate-origin labels are provisional and should not be equated with final eligibility decisions.}

\revise{The later assessment assigns topical relevance scores of 0 (no substantive connection), 1 (background), 2 (adjacent or transferable) and 3 (direct OS-security candidate). The 225 records contain 1, 37, 164 and 23 entries at these scores, respectively. Scores indicate topical fit rather than study quality or verified eligibility; the inventory records whether judgments used abstracts, accessible conclusions or earlier notes.}

\paragraph{\textbf{\revise{Screening personnel and limitations}}}
\revise{The surviving historical records do not document who screened each stage, independent duplicate screening or disagreement resolution. The later update does not establish independent human adjudication either. Inter-rater agreement cannot be reconstructed, and decisions near the core/adjacent boundary may be subjective.}

\paragraph{\textbf{\revise{Record counts and provenance}}}
\revise{Table~\ref{tab:selection-counts} distinguishes the historical scored-record inventory from the later merged candidate inventory. The archived historical file has 485 records, of which 220 carry scores 1--3: 45, 116 and 59 respectively. These counts match those in the earlier manuscript, which designated 59 score-3 records for in-depth review; the exact correspondence to the previously submitted supplementary file cannot be established from counts alone. Neither number is a verified count of unique eligible studies. The remaining 265 rows include scores of zero and missing or other values; they cannot be interpreted as documented exclusions. The later inventory draws on a different selection and updated ratings, so its totals cannot replace the historical numbers in a single sequential screening flow. A record-level crosswalk and unresolved matches are provided with the supplementary inventory.}

\begin{table}[htbp]
\centering
\caption{\revise{Available literature-record counts and their limits.}}
\label{tab:selection-counts}
\small
\begin{tabular}{@{}>{\RaggedRight\arraybackslash}p{0.44\columnwidth}>{\RaggedRight\arraybackslash}p{0.50\columnwidth}@{}}
\hline
\textbf{Quantity} & \textbf{Count/status} \\ \hline
\revise{Original search hits and selection stages} & \revise{Unknown; full search and decision logs unavailable} \\
\revise{Historical scored records} & \revise{220 records; 59 formerly designated for in-depth review} \\
\revise{Updated merged candidate inventory} & \revise{225 records, including 23 provisional score-3 records} \\
\revise{Separate discovery audit} & \revise{6,408 identifier groups; later arXiv pagination produced 1,168 additional unmatched identifiers} \\
\revise{Unique studies after version reconciliation} & \revise{Not established} \\
\revise{Unique eligible studies and final in-depth set} & \revise{Not established} \\ \hline
\end{tabular}
\par\smallskip
\parbox{0.96\columnwidth}{\footnotesize\revise{These counts describe different record sets, not consecutive stages. Preliminary scores do not certify eligibility.}}
\end{table}

\revise{The candidate and historical inventories and their interpretation are provided in \url{reproduct/}. Later-added provider reports are recorded separately in the literature addendum and are not included in those candidate counts. The accompanying Crossref metadata log documents the exploratory update; it does not reconstruct the earlier search. A separate citation-level audit maps manuscript references to the candidate and search records where identifiers or source-checked aliases permit; unmatched citations include background, adjacent work and unresolved cases, and are not automatically counted as omissions. Version aliases and source-supported corrections are retained with their evidence under \url{reproduct/audit/}, and the revised manuscript's citation mapping is in \url{reproduct/revision/}. Table~\ref{tab:llm-kernel-related} illustrates representative core and related work rather than a complete eligibility record. The KernelGPT workflow in the device-driver discussion is a literature-based illustration, not an experiment performed to validate the survey framework (\S~\ref{sec:os:driver}).}

\begin{table*}[t]
\caption{\revise{Representative works and their underlying techniques.}}
\label{tab:llm-kernel-related}
\centering
\footnotesize
\begin{tabular}{>{\RaggedRight\arraybackslash}p{0.22\textwidth}>{\RaggedRight\arraybackslash}p{0.72\textwidth}}
\hline
\textbf{\revise{Technique}} & \textbf{\revise{Representative work}} \\
\hline
\revise{Analysis and fuzzing} & \revise{KernelGPT refines syscall specifications; RandLuzz directs fuzzing on evaluated programs; SyzAgent guides kernel-test generation. Kernel applicability of RandLuzz requires separate evaluation}~\revcite{yang2025kernelgpt,feng2025fuzzing,syzagent25}. \\
\hline
\revise{Formal verification} & \revise{Selene studies seL4 proof generation, while AutoRocq and FVEL use prover feedback to refine proof attempts}~\revcite{zhang2024selene,tu2025agentic,lin2024fvel}. \\
\hline
\revise{Synthesis} & \revise{KNighter and Kgent generate checkers and eBPF extensions. QLCoder met its vulnerable-versus-patched detection criterion for 53.4\% of 176 CVEs across 111 Java projects; its applicability to kernel code requires separate evaluation}~\revcite{yang2025knighter,zheng2024kgent,wang2025qlcoder}. \\
\hline
\revise{Patching and repair} & \revise{kAgent and SAN2PATCH generate candidate repairs, while PortGPT targets kernel patch backporting; candidates require target-specific validation}~\revcite{mathai2025crashfixer,kim2025logs,li2025portgpt}. \\
\hline
\end{tabular}
\end{table*}
